\subsection{Book の低三角形：残る二つの sector と完全復元}
\label{nb20:section}
\label{nb21:section}

$G_{5,13}=\{01,02,03,04,12,13,14\}$ の二価 cores $2,3,4$ を扱う．
正の整数局所 charts，奇数高分母，低辺の node 非負・総和五，
対向低 $D$ 相殺の下で，canonical 非零低低 entry は $-L$，$L\in\{1,3,4\}$ となる．
全正枝，全負枝，二負枝は全て排除されており，残るなら負枝は一つ，正枝は二つである．
全ての低 $D$ は非零で，各 core は一正一負の低 $D$ を持つ．
これらは高中心 $0,1$ の完成を仮定しない必要条件である．

\paragraph{現在残る二つの sector．}
高 ray $0$ へ正 entry を持つ低 row の先を next とすると，有向 node 三角形になる．
canonical $A=0$ を $\sigma=+$，$A=1$ を $\sigma=-$ とする．
前者では next が $C$，後者では next が $D$ である．
next の低 $D$ の共通符号を $\tau$ とし，唯一の負枝 $N$ を $\sigma_N=+$ に正規化すると，
残る可能性は次だけである：
\[
 \mathrm I:\quad N^+\to P_1^+\to P_2^+\to N^+,
 \qquad \tau=+,\qquad (L_{P_1},L_{P_2})=(3,3),
\]
\[
 \mathrm {II}:\quad N^+\to P_1^+\to P_2^-\to N^+,
 \qquad \tau=+,\qquad (L_{P_1},L_{P_2})=(4,3).
\]
負枝は $S_N=T_N=w$，$H=r_N-w$ に対して
\[
 w\ge135,\qquad H\ge w
\]
を満たす．型 II ではさらに
\[
 S_{P_2}=w-1,\qquad T_{P_2}=w+1,
 \qquad n(P_2N)=(0,0,5)
\]
であり，node の順は高 ray $0,1$ と残る低 ray である．
正枝の小さい $L$ に対する符号・分母順序の補題と負枝の符号計算を合わせると，
この正規化の下で三つ全ての二価 core に
\[
 q_{i0}<q_{i1}\qquad(i=2,3,4)
\]
が成り立つ．これは高中心二つに接続する分母三つの厳密な比較である．
中心の四高分母条件への接続は未解決であり，book 自体の不存在とはしていない．

\paragraph{再利用する node offset．}
有向順序で next row を $(a_i,-b_i,-c_i)$，
previous row を $(-d_i,e_i,-f_i)$ と書く．最初の二成分は高 ray $0,1$ に対応する．
\[
 b_i=a_i+1-c_i,\quad d_i=e_i+1-f_i,\quad
 c_if_i=0,\quad c_i+f_i=L_i\in\{1,3,4\},
\]
\[
 (a_i,e_i,c_i,f_i)=
 \begin{cases}
 (r_i,T_i,0,L_i),&\sigma_i=+,\\
 (T_i,r_i,L_i,0),&\sigma_i=-.
 \end{cases}
\]
辺 $i\to i+1$ の node は
\[
 (n_0,n_1,n_{\mathrm{low}})
 =(1-a_i+d_{i+1},\ 1+b_i-e_{i+1},\ 1+c_i+f_{i+1})
\]
であり，非負性は有限 offset
\[
 c_i+f_{i+1}\le4,\qquad e_{i+1}=a_i+t_{i+1},\qquad
 f_{i+1}-2\le t_{i+1}\le2-c_i
\]
と同値である．残余 sector では $\tau=+$ なので，
負の $D$ 端から正の $D$ 端へ向けた周期は node 周期と逆向きになる．
二つの向きは定義が異なるため，計算中に同一視しない．
また高辺 $01$ があるので，負枝を禁じる低 wedge の補題は適用できない．

\paragraph{低係数と局所 sector の比較式．}
一般二高標準形の低係数は，$\delta=T-S$ として
\[
 D_C=\delta-\frac Sr s+\frac T{r+1}z,\qquad D_D=s-z-1.
\]
負枝では $w,H\ge2$，$w^2/2<H(H+1)<2w^2$ から
\[
 \frac{17}{10}w-\frac12<r<\frac{17}{7}w.
\]
$E_-=1-w^2/[H(H+1)]$ とすると，非零低 $D$ の下で
\[
 D_C>0\Longleftrightarrow E_->0\Longrightarrow r\ge2w,
 \qquad
 D_C<0\Longleftrightarrow E_-<0\Longrightarrow r\le2w-1.
\]
局所的に負と零の低 $D$ を持つ唯一の模型は
$(r,s,z,S,T)=(4,4,5,2,2)$，$(D_C,D_D)=(0,-2)$ であり，その低三角形完成は排除済みである．
零低辺を新たに残余 case として加え直す必要はない．

正枝では $\delta\in\{2,3\}$，$h=S-r$，$d=\delta-1$，$k=h+d$ と置く．
\[
 \alpha=\frac{sh}{r},\qquad\beta=\frac{zk}{r+1},\qquad
 F=(2r+1)h-r^2d,
\]
\[
 F\alpha\beta=(\beta-\alpha-1)hk,
 \qquad (rs+1)\beta-((r+1)z+1)\alpha=1.
\]
一正一負の低 $D$ を持つとき，$D_C$ と $F$ は同符号であり，
$D_C<0$ は $q_A>q_B$，$D_C>0$ は $q_A<q_B$ に対応する．
再利用する reverse bounds は
\[
 F>0\Longrightarrow
 T\ge\begin{cases}(3r+4)/2,&\delta=2,\\2r+3,&\delta=3,\end{cases}
 \qquad
 F<0\Longrightarrow
 T\le\begin{cases}\lfloor(3r+3)/2\rfloor,&\delta=2,\\2r+2,&\delta=3.\end{cases}
\]
$F<0$ では $-F\le hk$，従って
\[
 d r^2\le2hr+h(h+\delta).
\]
$\beta=1$ の reverse data は $(r,S,T,z)=(3,4,6,2),(2,3,6,1)$ だけである．
これは Bezout 恒等式から $k(1+2z-dz^2)=\delta$ を得て，
$k\ge\delta$ を使うと分かる．両共鳴族は node offset により負 $w\le5$ を強制するため，
現在の残余 sector では現れない．従って残余正枝の $F<0$ 側では $\beta\ge2$ であり，
\[
 0<\frac{-F}{hk}<\frac12,
 \quad 2d r^2<h^2+4rh+(\delta+1)h,
 \quad h>\frac37r-1,
 \quad T>\frac{10}{7}r+1
\]
を使える．

\paragraph{排除済み sector と有限 certificate の境界．}
二負枝を仮定すると，全ての $\sigma$ と $\tau$ の配置で負 core の一つに
$w\le50$ が導かれる．この解析的範囲の全 scale 復元は完成を排除する．
一負枝の二つの node 型では，上の局所不等式と有限 offset から
\begin{center}
\begin{tabular}{lc}
追加で排除した sector & 導出された負 $w$ の上界\\\hline
型 I，$\tau=-$ & $134$\\
型 I，$\tau=+$，少なくとも一方の正枝が $L=4$ & $8$\\
型 II，$\tau=-$ & $11$\\
型 II，$\tau=+$，$L_{P_1}=3$ & $6$
\end{tabular}
\end{center}
となる．$2\le w\le134$ の全 scale 復元では初期 reverse data が $6392$ 個，
負 canonical cores が $1089$ 個，左右の相手が合計 $169$ 個である．
両隣がともに非空なのは六 core だけで，最後の十一組の $D$ 残差は全て非零だった．
\nolinkurl{verification/verify_book_one_negative_sector_reduction_v21.py} と
対応する JSON が全初期 core・全相手集合・最後の十一組を保存している．
ここでは $w$ の上界自体が仮想的な sector から導かれている．
一方，残余二 sector の $w\ge135$ 側への一様上界はまだない．

\paragraph{固定 core の完全復元への入口．}
一つの二価 core の行列を固定すると，各低辺の二十一 node 分配が
相手の低 row 全体を $p^j_{ik}=1-n_k-p^i_{jk}$ に固定する．
対向 $D$ の相殺は，その低 ray の指定値 $D_*$ も固定する．
canonical reverse data $(r,S,T)$ は有限に復元できる：
固定 row が $C$ なら $r$ が固定され，正枝は $r<S<r\delta$，
負枝は $2\le w<r$ と既知の $H$ 範囲を満たす．
固定 row が $D$ なら $S,T$ が固定され，正枝は $S/\delta<r<S$，
負枝は $w^2/2<H(H+1)<2w^2$ を満たす．

$h=S-r$，$\delta=T-S\in\{0,2,3\}$ とすると，残る scale は
\[
 Fsz+As+Bz+C=0,\qquad \ell s+mz+n=0,
\]
\[
 F=(2r+1)h-r^2(\delta-1),\quad
 A=h(r+1),\quad B=-r(h+\delta-1),\quad C=r(r+1),
\]
\[
 (\ell,m,n)=
 \begin{cases}
 (-S(r+1),rT,(\delta-D_*)r(r+1)),&C\text{ ray},\\
 (1,-1,-D_*-1),&D\text{ ray}
 \end{cases}
\]
を満たす．$\gcd(2r+1,r^2)=1$ と $\delta\in\{0,2,3\}$ から $F\ne0$，
また $\ell m\ne0$ なので，消去した二次式
\[
 -F\ell s^2+(-Fn+Am-B\ell)s+(Cm-Bn)=0
\]
の最高次係数は非零である．固定 reverse data ごとに高々二つの scale を，
整数平方判定と整除から全て復元する．七・十三の共鳴族もこの二次式に含まれる．
二つの低隣接 core の有限集合を作り，残る低辺を検査すれば，全低三角形完成を得る．

負の固定 reverse data $(w,H)$ 自体の全 scale は
\[
 c=\frac{sH}{r},\quad e=\frac{z(H+1)}{r+1},\quad
 E_-=\frac MN\ne0,\quad E_-ce=e-c+1,
\]
\[
 (N-Mc)(N+Me)=N(N-M)>0
\]
から復元する．$c\ge2,e\ge1$ により両因子は正なので，全正約数を取ればよい．
この二つの復元法はいずれも $s,z$ または高分母の cutoff を必要としない．
汎用の固定 core 完成器は \nolinkurl{research/book_low_triangle_completion_v19.py} にあり，
別実装による第21版の整数行列・正の weights・相手復元は
\nolinkurl{verification/verify_book_two_negative_obstruction_v21.py} にある．
低 $D$ を元の pair equation に直接代入する後者は，旧 generator を import しない．

\paragraph{次に比較する量．}
残る二 sector に限り，node offset，低 $D$ の大きさ，
$q_{i0}<q_{i1}$ と中心二 core の四高分母条件を同時に扱う．
型 II の強制 row $S_{P_2}=w-1,T_{P_2}=w+1$ は逆復元の入口になる．
固定 core の有限性を core 自体が動く族の有限性と取り違えず，
新しい一様上界または反例構成を優先する．
